\documentclass[11pt]{article}

\usepackage[T1]{fontenc}
\usepackage[utf8]{inputenc}
\usepackage{lmodern}
\usepackage{amsmath,amssymb,amsthm,bm}
\usepackage{booktabs}
\usepackage{graphicx}
\usepackage{microtype}
\usepackage{hyperref}
\usepackage[margin=1in]{geometry}

\newtheorem{proposition}{Proposition}
\newcommand{\Tr}{\operatorname{Tr}}
\newcommand{\dd}{\mathrm{d}}
\newcommand{\RQIR}{\mathrm{RQIR}}

\title{Relativity--Quantum Interface Reconstruction I:\\
Operational Hierarchy, Ordered Stress--Energy Information, and Finite Calibration-Nullspace Discriminants}
\author{Authors to be fixed before submission}
\date{Working manuscript v0.5 / 8 September 2026}

\begin{document}
\maketitle

\begin{abstract}
Laboratory tests at the gravity--quantum interface are often formulated either by selecting a candidate microscopic theory or by proposing a particular witness of gravitational nonclassicality. We develop a complementary operational reconstruction problem: which components of quantum source information remain distinct after a declared finite calibration, before a microscopic interpretation of gravity is assigned? Building on the established closed-time-path and stochastic-gravity distinction between one-point stress--energy, symmetrized fluctuations, and causal response, we organize the source information into an ordered hierarchy containing $J=\langle T\rangle$, the symmetrized noise kernel $N$, the antisymmetric/commutator sector $D$, the retarded response $\chi^R$, and higher closed-time-path information. We then treat finite source calibration as a quotient on the physical state space. A finite-dimensional proposition shows that if the calibration map has a physical null direction on which a specified ordered-response functional is nonzero, then two positive source states can agree on every declared calibration observable while differing in that response. Constructive RQIR examples realize positivity, rank, calibration equality and detector-facing response simultaneously, and show that changing calibration geometry rotates the hidden direction. The result is an operational source-level distinguishability statement, not evidence that gravity is quantized and not a claim that gravity necessarily couples to $D$ or $\chi^R$. Classical, semiclassical, stochastic, measurement-feedback, postquantum and quantum-gravity descriptions can overlap in observable space and must be compared through the complete interface map. Paper I therefore supplies a calibration-aware source-information layer for the nuisance-profiled statistical analysis developed in Paper II.
\end{abstract}

\section{Introduction}

The empirical interface between gravitation and quantum theory is increasingly accessible to laboratory analysis. Proposed probes include matter-wave and spin interferometry, gravitationally mediated correlations and entanglement, quantum clocks \cite{Roura2020,Wakakuwa2026}, optomechanical tests, decoherence and diffusion measurements, field-disturbance tests, and observables designed to go beyond a static Newtonian-potential description. These programmes differ in experimental architecture and in what they aim to infer, but they share a basic inverse problem: the observed probability distribution is generated by a chain of source preparation, quantum-matter dynamics, gravitational or effective interface dynamics, calibration, detector transfer, nuisance parameters and data reduction.

Consequently, a nonclassical-looking signal need not by itself determine a unique microscopic description of gravity. Conversely, microscopic models that are conceptually different may be operationally degenerate in a finite experimental domain. The appropriate inference target is therefore not automatically a theory label but an equivalence class of interface maps compatible with the declared observables and controls.

Relativity--Quantum Interface Reconstruction (RQIR) begins from this inverse perspective. Schematically,
\begin{equation}
  P_{\rm data}(\bm{o}\mid\bm{s}) \longrightarrow [\mathfrak I],
\end{equation}
where $[\mathfrak I]$ denotes an operational equivalence class of gravity--quantum interface maps after declared calibrations, baselines and nuisance freedoms are taken into account. RQIR does not begin by assuming that gravity is fundamentally classical, stochastic, quantum, emergent or hybrid.

Paper I isolates the earliest source-side layer of this programme. Its central question is narrower than ``is gravity quantum?'':
\begin{quote}
\emph{What source information can remain operationally distinct after a declared finite calibration map?}
\end{quote}

The analysis uses established quantum-field-theoretic and open-system structures. Stochastic gravity and influence-functional methods already distinguish the stress--energy expectation value from its symmetrized fluctuation kernel and from response/dissipation structure \cite{FeynmanVernon1963,Keldysh1965,Kubo1957,HuVerdaguer2008}. RQIR does not claim this separation as new. Instead, we use these objects as coordinates of an operational discrimination problem and ask what a finite calibration can and cannot identify.

This leads to a geometric statement. If a finite calibration map leaves a physical null direction and a specified ordered-response functional does not vanish on that direction, then two positive source states can be calibration-indistinguishable while remaining response-distinct. The proposition is elementary at the level of linear algebra; its role is to make the calibration quotient explicit in the gravity--quantum source problem and connect it to constructive positive-state examples and detector-facing response.

The interpretation is deliberately conservative. Paper I does not assert that gravity couples to the matter commutator or retarded stress--energy response in nature. It does not assert that a nonzero matter commutator proves quantized geometry. It does not treat observation of entanglement as an assumption-free proof of quantum gravity. Instead, it separates source-side information structure from the later questions of statistical identifiability, physical resources and model/comparator uniqueness.

\subsection{Contributions}

Within the frozen Paper-I scientific scope, the contribution is the combined structure
\begin{enumerate}
  \item an operational reconstruction framing in which the inferred object is an interface equivalence class rather than a presupposed microscopic gravity theory;
  \item an ordered source-information hierarchy using established $J$, $N$, $D$, $\chi^R$ and higher CTP objects as distinct operational coordinates unless a physical identity relates them;
  \item a finite calibration quotient defined on physically admissible source-state perturbations;
  \item a calibration-nullspace proposition identifying when positive calibration-equivalent states can remain distinct under a specified ordered-response functional;
  \item constructive RQIR realizations verifying positivity, rank, calibration equality and response separation together;
  \item a calibration-steering relation showing that the hidden direction depends on calibration geometry and is therefore a design variable;
  \item an explicit epistemic boundary between source-side distinguishability and claims about the ontology of gravity;
  \item a controlled bridge from exact source separation to nuisance-profiled statistical identifiability in Paper II.
\end{enumerate}

\subsection{Non-claims}

Paper I does not claim empirical evidence for new physics or quantized geometry; that gravity necessarily transmits $D$ or $\chi^R$; that the mean/noise/response split is an RQIR discovery; that the generic incompleteness of finite quantum measurements is new; that positivity constraints in quantum tomography are new; that quotient geometry for calibration or nuisance uncertainty is new; that entanglement alone is sufficient under every admissible comparator class; relativistic completeness of the finite toy models; detector-level identifiability after nuisance profiling; or apparatus-specific resource closure.

\section{Relation to Established Approaches}

\subsection{Semiclassical and stochastic gravity}

Standard semiclassical gravity uses the renormalized stress--energy expectation value as the source of a classical metric. Schematically,
\begin{equation}
 G_{\mu\nu}+\Lambda g_{\mu\nu}+H^{\rm EFT}_{\mu\nu}
 = 8\pi G\,\langle \hat T_{\mu\nu}\rangle_{\rm ren}.
\end{equation}
Stochastic gravity extends the mean-field description by including stress--energy fluctuations through the noise kernel and, in influence-functional/CTP formulations, contains fluctuation and response/dissipation structure \cite{HuVerdaguer2008}. Thus a fluctuation-sensitive signal is not automatically an intrinsically quantum-gravitational signal.

\subsection{Information-theoretic witnesses}

Gravity-mediated entanglement and related information-theoretic witnesses provide one major route to laboratory tests \cite{BoseEtAl2017,MarlettoVedral2017,MarlettoVedral2020,MarlettoVedral2025RMP}. Non-entanglement LOCC-based inequalities provide another route \cite{LamiPedernalesPlenio2024}. RQIR treats these as operational constraints whose interpretation depends on the admitted comparator class rather than privileging one witness as the definition of gravitational quantumness.

\subsection{Classical, hybrid and postquantum alternatives}

Classical--quantum coupling models enlarge the relevant comparator space. Postquantum classical gravity, irreversibility constraints, stochastic decoherence/diffusion constructions and measurement-feedback or conditional semiclassical descriptions illustrate this point \cite{Oppenheim2023,GalleyGiacominiSelby2023,OppenheimEtAl2023,Diosi2025,MikiEtAl2025}.

\subsection{Informational incompleteness, positivity and calibration quotients}

Informationally incomplete quantum tomography explicitly studies non-uniqueness of density operators compatible with finite measurement records \cite{DArianoParisSacchi2003,TeoRehacekHradil2013,GoncalvesEtAl2013}. Positivity can strongly restrict, and for suitable sensing maps even remove, otherwise available reconstruction ambiguity \cite{KalevKosutDeutsch2015}. Conversely, an informationally incomplete record can still be conclusive for a specified membership question without reconstructing the entire state \cite{LyyraKuuselaHeinosaari2019}. Paper I does not claim any of these general facts as new.

There is also a distinct literature on kernels and invertibility of retarded response maps. In density-functional-like settings, Giesbertz characterized conditions under which perturbing potentials lie in the kernel of a retarded response function, first for ground-state and then for mixed-state settings \cite{Giesbertz2015,Giesbertz2016}. That kernel is the kernel of a dynamical response map from perturbations to observables. The nullspace in RQIR-THM-001 is different: it is the unresolved direction of a declared finite calibration map acting on physically admissible source-state perturbations, and the retarded or ordered response is evaluated as a separate discriminant on that calibration nullspace.

Quotient geometry under calibration uncertainty has also appeared recently in quantum metrology. Wani and Al-Kuwari profile nuisance directions in the quantum Fisher metric and define a quotient quantum-speed limit \cite{WaniAlKuwari2026}. This is conceptually aligned with the general principle that calibration ambiguity should be quotiented before inference, but the object being quotiented differs from Paper I: their construction profiles parameter-estimation directions in a state-space information metric, whereas RQIR begins from a finite source calibration map and asks whether a designated ordered-response functional survives its physical nullspace. The nuisance-profiled information geometry is closer to the statistical layer developed in Paper II.

\subsection{Current witness debate and dynamical probes}

Aziz and Howl argued that local classical gravity with matter treated in full QFT can produce entanglement through matter-mediated channels \cite{AzizHowl2025}. Subsequent analyses dispute the relevance or admissibility of some proposed counterexamples in particular regimes; for example, Feng, Vedral and Marletto analyze collapse-based models and locality assumptions \cite{FengVedralMarletto2026}. We do not attempt to settle this active debate. Its methodological implication is sufficient here: witness interpretation is conditional on the comparator class and assumptions defining the full observable map.

Dynamical field-sensitive probes can also access information beyond a static Newton-potential description \cite{ChenGiacomini2025}. This motivates an operational hierarchy broad enough to retain causal response and higher ordered information.

\section{Operational Reconstruction of the Interface}

Let $\mathcal O=\{O_A\}$ denote operational observables, where $A$ includes state preparation, smearing/coarse graining, detector transfer, experimental settings and any renormalization prescription required for the observable to be meaningful. For an explicit baseline $B$ in domain $\mathcal D$, define
\begin{equation}
 \Delta_A^{(B)}(\mathcal D)
 = O_A^{\rm obs}(\mathcal D)-O_A^{(B)}(\mathcal D).
\end{equation}
Residuals relative to distinct baselines are not combined without an explicit transformation.

Abstractly,
\begin{equation}
 \mathfrak I:(\rho_{\rm matter},\mathcal G,\mathcal C,\lambda)
 \longrightarrow P(\bm{o}\mid\bm{s}),
\end{equation}
where $\mathcal G$ denotes gravitational/interface degrees of freedom or effective variables, $\mathcal C$ denotes constraints and consistency structure, and $\lambda$ collects model and nuisance parameters.

\section{Ordered Stress--Energy Information}

Define
\begin{equation}
 \delta\hat T_A(x)=\hat T_A(x)-\langle\hat T_A(x)\rangle.
\end{equation}
The one-point source is
\begin{equation}
 J_A(x)=\langle\hat T_A(x)\rangle.
\end{equation}
The greater and lesser correlators are
\begin{align}
 G^>_{AB}(x,y)&=\langle\delta\hat T_A(x)\delta\hat T_B(y)\rangle,\\
 G^<_{AB}(x,y)&=\langle\delta\hat T_B(y)\delta\hat T_A(x)\rangle.
\end{align}
Define the symmetrized/noise coordinate
\begin{equation}
 N_{AB}(x,y)=\frac12\langle\{\delta\hat T_A(x),\delta\hat T_B(y)\}\rangle
 =\frac12(G^>+G^<),
\end{equation}
and the antisymmetric/commutator coordinate
\begin{equation}
 D_{AB}(x,y)=\frac{1}{2i}\langle[\delta\hat T_A(x),\delta\hat T_B(y)]\rangle,
 \qquad G^>-G^<=2iD.
\end{equation}
With a declared sign convention, the retarded response is
\begin{equation}
 \chi^R_{AB}(x,y)=\frac{i}{\hbar}\theta(x^0-y^0)
 \langle[\hat T_A(x),\hat T_B(y)]\rangle.
\end{equation}

The central ordering-preservation rule is that matching $J$ and $N$ does not imply matching $D$ or $\chi^R$ unless an additional physical identity, commutativity condition, spacelike-separation statement, equilibrium relation or controlled limit supplies that implication. This is compatible with fluctuation--dissipation relations: those relations are additional physical structure valid in specified regimes, not a universal algebraic identity between arbitrary calibrated states \cite{Kubo1957}.

\section{Closed-Time-Path Parent Object}

A compact parent object is the Schwinger--Keldysh/closed-time-path generating functional
\begin{equation}
 Z_T[J_+,J_-]=\Tr\!\left(U[J_+]\rho_T U[J_-]^\dagger\right),
 \qquad W_T=-i\hbar\ln Z_T.
\end{equation}
Functional derivatives generate branch-ordered correlators, including Wightman, time-ordered, retarded and higher nested-response structures \cite{FeynmanVernon1963,Keldysh1965}. RQIR uses this established formalism as an organizational parent object rather than claiming it as new.

\section{Finite Calibration as an Operational Quotient}

Let $V$ be the real tangent space of Hermitian source-state perturbations after exact hard linear preparation constraints have been eliminated. Let
\begin{equation}
 A:V\rightarrow\mathbb R^m
\end{equation}
be a finite calibration map, and let
\begin{equation}
 c:V\rightarrow\mathbb R
\end{equation}
be a specified ordered-response functional. If $\ker A$ is nontrivial, calibration leaves an equivalence class. If $c$ vanishes on the entire nullspace, that response is fixed by calibration. If a physical $n\in\ker A$ satisfies $c(n)\neq0$, the response remains distinct within the calibration-equivalence class.

This geometry is related to informational incompleteness in quantum tomography \cite{DArianoParisSacchi2003,TeoRehacekHradil2013,GoncalvesEtAl2013,KalevKosutDeutsch2015}. The RQIR-specific use is to impose physical-state admissibility and evaluate a designated ordered-response coordinate on the unresolved direction, rather than to reconstruct a unique global state estimator.

\section{Finite Calibration-Nullspace Response Discriminant}

\begin{proposition}[RQIR-THM-001]\label{prop:rqir-thm-001}
Let $V$ be the real tangent space of Hermitian source-state perturbations after all exact hard linear constraints in the declared preparation domain have been eliminated. Let $A:V\to\mathbb R^m$ be a finite linear calibration map with one-dimensional nullspace $\ker A=\mathrm{span}\{n\}$, and let $c:V\to\mathbb R$ be a linear response functional. Suppose that $c(n)\neq0$, the nominal density operator $\rho_0$ lies in the interior of the physical state set on the retained finite-dimensional Hilbert space, and $n$ is represented by a Hermitian traceless perturbation compatible with the exact preparation constraints. Then there exists $\epsilon>0$ such that
\begin{equation}
 \rho_\pm=\rho_0\pm\epsilon n
\end{equation}
are physical states, have identical declared calibration data, and differ in response.
\end{proposition}

\begin{proof}
Because $\rho_0$ is strictly positive, choose
\begin{equation}
 0<\epsilon<\frac{\lambda_{\min}(\rho_0)}{\lVert n\rVert_{\rm op}}.
\end{equation}
Weyl's eigenvalue bound guarantees positivity of $\rho_0\pm\epsilon n$. Trace and the exact linear preparation constraints are preserved because $n$ lies in the reduced tangent space. Since $An=0$, the declared calibration rows agree for the pair. Since $c(n)\neq0$,
\begin{equation}
 c(\rho_+-\rho_-)=2\epsilon c(n)\neq0.
\end{equation}
\end{proof}

The generic statement that incomplete finite measurements can possess null directions is established measurement theory. The purpose of RQIR-THM-001 is narrower: it packages a positivity-preserving source-state pair inside a declared calibration quotient and evaluates a designated ordered-response functional on that physical direction. Likewise, RQIR-THM-001 is not an invertibility theorem for the retarded response operator itself; the response functional $c$ is the discriminator evaluated after the calibration quotient has been defined.

\section{Constructive Realization and Calibration Steering}

The deterministic constructions Toy009 and Toy010 test the hypotheses of Proposition~\ref{prop:rqir-thm-001}; they do not replace its proof. Both use a 25-dimensional Hermitian source parameterization and realize the codimension-one regime $\operatorname{rank}A=24$ under the declared numerical-rank rule. A clean independent GitHub Actions rerun on 8 September 2026, with Python 3.12.14 and NumPy 2.1.3, reproduced the frozen Toy009/Toy010 regression checks at repository commit \texttt{763757788397564d180cfd6b33c28db57412d74c}.

The constructive certificate has six logically separate checks:
\begin{enumerate}
 \item establish the codimension-one rank condition under a stated numerical-rank rule;
 \item establish a one-dimensional numerical nullspace and record a normalized hidden direction $n$;
 \item verify physical-state admissibility, including positivity of the paired states;
 \item verify equality of the declared calibration coordinates, including the selected mean/noise coordinates;
 \item evaluate the ordered and detector-aware response coordinates on the same frozen direction and verify a nonzero separation under a stated decision rule;
 \item bind every quoted numerical value to the immutable run artifact, exact code revision, configuration and precision policy that generated it.
\end{enumerate}
Items 1--5 are mathematical or numerical checks of the construction. Item 6 is a provenance requirement for manuscript-level numerical claims. All six checks are now satisfied by the frozen clean-run certificate. In Toy009 the maximum equality residual is $2.2\times10^{-16}$, the selected mean/noise differences are below $2\times10^{-18}$, both paired density operators are positive, and the ordered responses are opposite with magnitude approximately $1.21\times10^{-2}$. Keeping the Toy009 source fixed while changing only the finite calibration geometry, Toy010 retains $\operatorname{rank}A=24$, gives a maximum equality residual $4.4\times10^{-16}$, retains positive paired states and equal selected mean/noise coordinates, and produces opposite ordered responses of magnitude approximately $1.33\times10^{-2}$. The normalized null directions of the inherited Toy009 and Toy010 calibrations differ by approximately $37.7^\circ$. These values are finite-dimensional construction diagnostics rather than experimental forecasts or evidence that gravity transmits the corresponding response coordinate.

Calibration is also an active design variable. More explicitly, let $A(t)\in\mathbb R^{(p-1)\times p}$ be differentiable and of constant full row rank in a neighborhood, and choose a differentiable normalized null vector $n(t)$ satisfying
\begin{equation}
 A(t)n(t)=0,\qquad \lVert n(t)\rVert_2=1.
\end{equation}
Differentiating the null relation gives
\begin{equation}
 A'n+An'=0.
\end{equation}
Since $AA^+=I_{p-1}$, the general local solution can be written
\begin{equation}
 n'=-A^+A'n+\alpha n.
\end{equation}
For full row rank, $\operatorname{ran}(A^+)$ is the row space of $A$ and is orthogonal to $\ker A$. Differentiating the normalization condition gives $n^Tn'=0$, hence $\alpha=0$ and therefore
\begin{equation}
 n'=-A^+A'n.
\end{equation}
In the operator 2-norm,
\begin{equation}
 \lVert n'\rVert_2\leq \lVert A^+\rVert_2\lVert A'\rVert_2
 =\frac{\lVert A'\rVert_2}{\sigma_{\min}(A)},
\end{equation}
where $\sigma_{\min}(A)$ is the smallest nonzero singular value. This is a local constant-rank conditioning statement; it is not a guarantee through a singular-value crossing. If the response functional $c$ is fixed while only calibration geometry changes, then
\begin{equation}
 \frac{\dd}{\dd t}c(n)=-cA^+A'n.
\end{equation}
Thus calibration geometry can rotate the unresolved direction and alter the response coordinate without changing the underlying source Hamiltonian.

\subsection{Evidence classes and promotion rule}

To prevent an implementation diagnostic from silently acquiring the status of an externally validated numerical claim, Paper I separates three evidence classes:
\begin{table}[ht]
\centering
\small
\begin{tabular}{@{}p{0.24\linewidth}p{0.45\linewidth}p{0.22\linewidth}@{}}
\toprule
Evidence class & Minimum frozen record & Manuscript status \\
\midrule
Analytic existence statement & Explicit hypotheses and proof & Publication claim \\
Constructive toy witness & Code revision, configuration, rank/null diagnostics, physical-state checks, calibration-equality record and response evaluation & Publication evidence only after immutable artifact binding \\
External numerical comparison & Source authority, convention, extraction rule and complete run provenance & Withheld until authority is frozen \\
\bottomrule
\end{tabular}
\caption{Promotion rule separating analytic claims, internally reproducible constructive evidence, and externally anchored numerical claims.}
\label{tab:evidence-promotion}
\end{table}

This rule makes the submission boundary operational. Missing run provenance can delay promotion of a toy number without weakening Proposition~\ref{prop:rqir-thm-001}; missing external authority can delay a calibrated comparison without changing either the analytic theorem or an internally reproduced constructive witness.

\section{Operational Sensitivity Hierarchy}

RQIR uses the working hierarchy
\begin{description}
 \item[L1 mean-sensitive:] only $J=\langle T\rangle$ is operationally relevant;
 \item[L2 noise-sensitive:] $N$ contributes independent information;
 \item[L3 order/response-sensitive:] $D$ or $\chi^R$ contributes independently;
 \item[L4 higher-order-sensitive:] higher CTP cumulants or nested responses are required;
 \item[L5 quantum-information-sensitive:] entanglement, channel or process observables add independent constraints.
\end{description}
These levels classify transmitted operational information, not microscopic gravitational ontology.

\section{Degeneracy Discipline}

Agreement with $\langle T_{\mu\nu}\rangle$ tests only a mean-source coordinate. A fluctuation excess is not automatically an intrinsic quantum-gravity fluctuation because stochastic gravity, matter-induced fluctuations, detector noise and environmental noise can populate overlapping observables. A nonzero matter commutator or source response is not by itself evidence for an operator-valued metric. Likewise, entanglement or another nonclassical correlation can be a powerful witness only relative to specified assumptions and comparator classes \cite{MarlettoVedral2025RMP,AzizHowl2025,FengVedralMarletto2026}.

Low-energy gravitational EFT supplies a controlled quantum comparator without choosing a UV completion \cite{Donoghue1994}. It is therefore a later RQIR comparator, not automatic evidence that a new theory is required.

\section{Structural-Identifiability Bridge}

Later RQIR comparator work exposed a distinct obstruction that can be stated without reopening the Paper-I scope. Exact on-shell information need not uniquely determine an off-shell/source-completed response object. Schematically,
\begin{equation}
 \Gamma_B=\Gamma_A+K_2(p)H,\qquad K_2(p)=0\ \text{on shell},
\end{equation}
can leave on-shell data unchanged while an off-shell/source-completed derivative differs. The Paper-I lesson is structural: a compressed observable map can quotient information needed by a source-completed response. Detailed comparator/source/Ward/IR calculations belong to Paper IV.

\section{From Exact Separation to Statistical Identifiability}

Paper I asks whether a declared calibration leaves a response-sensitive physical direction,
\begin{equation}
 An=0,\qquad c(n)\neq0.
\end{equation}
This exact statement does not establish that a finite noisy experiment can infer the associated amplitude after source, apparatus and detector uncertainties are introduced. Paper II addresses the nuisance-profiled information,
\begin{equation}
 F_{\beta\mid\theta}=F_{\beta\beta}
 -F_{\beta\theta}F_{\theta\theta}^{-1}F_{\theta\beta},
\end{equation}
or, in whitened geometry,
\begin{equation}
 F_{\beta\mid\theta}=\lVert (I-P_J)\tilde s\rVert^2.
\end{equation}
Thus exact source-side distinguishability is necessary but not sufficient for experimental inferability. Calibration-aware quotient information geometry in quantum metrology \cite{WaniAlKuwari2026} provides a useful neighboring example of why nuisance directions must be removed before an operational bound is interpreted.

\section{Discussion}

Paper I moves the gravity--quantum interface problem one level upstream from a binary choice of microscopic ontology. The ordered source hierarchy is established physics; the RQIR step is to make it part of an inverse operational map. Calibration then defines an explicit quotient on physical source states. RQIR-THM-001 records when that quotient leaves a positivity-compatible response-sensitive direction, while Toy009/Toy010 provide constructive certificates. The calibration-steering relation further promotes quotient geometry from passive bookkeeping to an experimental co-design variable.

The novelty claim is intentionally conjunctive rather than elemental. Informational incompleteness, positivity-constrained state recovery, response-function kernels and quotient geometry under nuisance/calibration uncertainty all have established antecedents \cite{TeoRehacekHradil2013,KalevKosutDeutsch2015,Giesbertz2015,Giesbertz2016,WaniAlKuwari2026}. The Paper-I contribution is the calibration-aware combination in which a finite source calibration map is imposed on physically admissible state perturbations, a positive pair is retained inside one calibration-equivalence class, a separate ordered-response functional is evaluated on the unresolved direction, and calibration geometry is treated as a controllable steering variable. This is the narrow claim to be tested against further prior-art searches before submission.

The epistemic staging is essential. Paper I identifies exact source-information directions. Paper II asks whether those directions survive nuisance profiling. Paper III converts surviving directions into physical resource and architecture requirements. Paper IV compares existing gravity and quantum-gravity realizations through a common funnel. Paper V is conditional and is justified only if that comparator analysis leaves a robust residual requiring a genuinely new parent model.

\section{Conclusion}

We formulated the first layer of RQIR as a calibration-aware source-information problem. Established closed-time-path and stochastic-gravity structures provide distinct mean, noise and response coordinates. A finite calibration map defines an operational equivalence class on physical source states. If that map leaves a physical null direction on which a specified ordered-response functional is nonzero, then positive calibration-equivalent but response-distinct states exist. Constructive RQIR examples realize the rank, positivity, equality and detector-facing response requirements together, and calibration geometry can rotate the hidden direction.

The result does not establish that gravity transmits the response coordinate, does not demonstrate quantized geometry and does not select a unique microscopic theory. It supplies a source-side operational structure to be passed to statistical, resource and comparator stages.

\section*{Planned Figures and Tables}

\begin{enumerate}
 \item RQIR inverse-problem architecture.
 \item Ordered second-order decomposition $G^>=N+iD$, $G^<=N-iD$ and causal projection to $\chi^R$.
 \item Physical-state/calibration-quotient geometry for RQIR-THM-001.
 \item Calibration steering under a perturbed calibration map.
 \item Epistemic ladder: exact separation $\to$ statistical identifiability $\to$ resource closure $\to$ comparator uniqueness $\to$ theory interpretation.
\end{enumerate}

Planned tables: operational sensitivity levels; literature positioning; reviewer-facing claims/non-claims boundary; evidence-class promotion rule in Table~\ref{tab:evidence-promotion}.

\section*{Reproducibility Note}

The Toy009/Toy010 numerical statements are bound to the frozen generation scripts and the clean-run authority recorded in \texttt{manuscripts/paper\_I/TOY009\_TOY010\_REPRO\_MANIFEST.md}. Workflow run \texttt{34177686395} completed successfully at commit \texttt{763757788397564d180cfd6b33c28db57412d74c}; the uploaded certificate artifact is \texttt{paper-i-toy009-toy010-certificate}, artifact ID \texttt{10037815979}, with archive digest \texttt{sha256:6363cfe9f5fad3ad7d04363921cc6d66e051ce7beaf9e22a3940a1b38c654b95}. The record includes the exact code revision, seeds/configuration, Python/NumPy environment, numerical-rank policy, positivity/equality diagnostics, ordered-response values and file hashes. A long-retention textual mirror of the clean output is stored at \texttt{manuscripts/paper\_I/certificates/TOY009\_TOY010\_CLEANRUN\_2026-09-08.md}, so the manuscript authority does not depend permanently on the temporary Actions retention window.

Any quantity calibrated to or compared with an external source additionally requires the authority record, convention and extraction rule for that source. A manuscript-level external numerical statement is promotable only when all fields relevant to its evidence class in Table~\ref{tab:evidence-promotion} are frozen. Missing external authority does not invalidate the analytic theorem or the internally reproduced Toy009/Toy010 certificate; it blocks only promotion of the affected external comparison or calibrated number.

\bibliographystyle{unsrt}
\bibliography{../../docs/PAPER_I_REFERENCES_WORKING}

\end{document}
